% !TeX root = ../main.tex
\chapter{Introduction}
\label{ch:intro}

\section{Purpose}

This document records the architectural choices behind the GENESYS Translation
Tool and the reasoning for each. It answers four questions: which architectural
patterns the system uses, how its subsystems fit together, whether those choices
follow from the original requirements, and which object-oriented design patterns
the implementation will leverage.

The System Requirements Specification
states what the tool must do and assigns every requirement an identifier. The
UML Design Document models the structure and behaviour. The Tool Selection
Document records which external tools and interfaces the project builds on and
why the alternatives were set aside. This document explains why the structure
takes the shape it does. Requirement identifiers are those of
the requirements specification and are used here unchanged, so that a claim made
about the architecture can be checked against the requirement that motivated it.

\section{The System in Brief}

The GENESYS Translation Tool moves behavioural design between a GENESYS
repository and a source code repository, in both directions, and leaves the two
reconcilable afterwards. The activity diagram is the pivot, and translation is
specified as two reversible pairs that meet there: activity diagram to and from
state diagram, both ends inside GENESYS, and activity diagram to and from code.

The governing criterion, taken from the client, is that an engineer should be
able to read a generated diagram and understand the system's behaviour without
reading its code. It is stated as \req{NFR-8}. Where it conflicts with
structural fidelity to the source, readability takes priority. That criterion
shapes the architecture directly, because it is the reason an Abstraction
subsystem exists at all.

\section{Document Conventions}

Every figure carries a number, a caption and a label, is referenced from the
surrounding text, and is followed by a table or by prose explaining its
elements. Named architectural patterns are set in small capitals, as
\pattern{Pipe and Filter}, so that a pattern claim is distinguishable from
ordinary use of the same words. The sources these names are taken from are
listed in Chapter~\ref{ch:references}.

\begin{itemize}[nosep]
	\item \req{CM-n} --- common to all translations.
	\item \req{SA-n} --- state diagram to activity diagram.
	\item \req{AS-n} --- activity diagram to state diagram.
	\item \req{AC-n} --- activity diagram to code.
	\item \req{CA-n} --- code to activity diagram.
	\item \req{FS-n} --- stretch functional requirements.
	\item \req{NFR-n} --- non-functional, minimum viable product.
\end{itemize}

Constraints are \textbf{C-n} and assumptions \textbf{A-n}. A constraint is
something the environment imposes; a design decision taken freely is never
recorded as one. The distinction matters here, because several of the choices
below were not choices at all.

\section{How This Document Is Organised}

\begin{longtable}{L{0.30\textwidth}L{0.62\textwidth}}
	\caption{Chapters against the questions they answer}\label{tab:organisation}\\
	\toprule
	\textbf{Chapter} & \textbf{Question answered} \\
	\midrule
	\endfirsthead
	\toprule
	\textbf{Chapter} & \textbf{Question answered} \\
	\midrule
	\endhead

	Chapter~\ref{ch:patterns}, Architectural Patterns & Which architectural
	patterns are used, which were considered and set aside, and why. Every
	candidate is given a verdict, not only the one adopted. \\
	\addlinespace
	Chapter~\ref{ch:subsystems}, Subsystems & What the subsystems are, which
	process each belongs to, and how they work together to produce a
	translation. \\
	\addlinespace
	Chapter~\ref{ch:design-patterns}, Design Patterns & Which object-oriented
	design patterns the implementation leverages, and what each one is there
	to solve. \\
	\addlinespace
	Chapter~\ref{ch:requirements}, Requirements Behind the Architecture &
	Whether these choices follow from the original requirements, traced
	requirement by requirement. \\
	\bottomrule
\end{longtable}
